\documentclass[10pt, aspectratio=169]{beamer}
% Preámbulo modular para presentaciones Beamer (Modelación Estadística)
% Diseño y paleta institucional del Tecnológico de Monterrey

\documentclass[aspectratio=169,xcolor={dvipsnames,table}]{beamer}

\usepackage[utf8]{inputenc}
\usepackage[spanish,mexico]{babel}
\usepackage{amsmath,amssymb,amsfonts,mathtools}
\usepackage{graphicx}
\usepackage{listings}
\usepackage{booktabs}
\usepackage{tikz}
\usetikzlibrary{matrix,arrows,decorations.pathmorphing,shapes.geometric,positioning}

% Paleta Institucional Tecnológico de Monterrey
\definecolor{TecAzul}{HTML}{0039A6}       % Azul TEC: color primario institucional
\definecolor{TecAzulOscuro}{HTML}{1A2E51} % Azul marino oscuro
\definecolor{TecRojo}{HTML}{EC2661}       % Rojo de acento (paleta miTec)
\definecolor{TecGris}{HTML}{646464}       % Gris neutro para comentarios y subtítulos
\definecolor{TecFondoBloque}{HTML}{F4F6F9}% Fondo suave para bloques

% Configuración de Tema Beamer
\usetheme{Boadilla}
\usecolortheme[named=TecAzulOscuro]{structure}
\setbeamercolor{alerted text}{fg=TecRojo}
\setbeamercolor{palette primary}{bg=TecAzulOscuro,fg=white}
\setbeamercolor{palette secondary}{bg=TecAzul,fg=white}
\setbeamercolor{palette tertiary}{bg=TecAzulOscuro!80!black,fg=white}
\setbeamercolor{title}{fg=TecAzulOscuro,bg=TecFondoBloque}
\setbeamercolor{frametitle}{fg=TecAzulOscuro,bg=TecFondoBloque}

% Bloques personalizados elegantes
\setbeamercolor{block title}{bg=TecAzulOscuro,fg=white}
\setbeamercolor{block body}{bg=TecFondoBloque,fg=black}
\setbeamercolor{block title alerted}{bg=TecRojo,fg=white}
\setbeamercolor{block body alerted}{bg=TecRojo!8,fg=black}
\setbeamercolor{block title example}{bg=TecAzul,fg=white}
\setbeamercolor{block body example}{bg=TecAzul!8,fg=black}

% Redondear bordes y añadir sombra sutil
\setbeamertemplate{blocks}[rounded][shadow=true]
\setbeamertemplate{navigation symbols}{}

% Pie de página limpio con número de diapositiva
\setbeamertemplate{footline}{
  \leavevmode%
  \hbox{%
  \begin{beamercolorbox}[wd=.7\paperwidth,ht=2.25ex,dp=1ex,left]{author in head/foot}%
    \hspace*{2ex}\insertshortauthor~~(\insertshortinstitute) \hfill \insertshorttitle
  \end{beamercolorbox}%
  \begin{beamercolorbox}[wd=.3\paperwidth,ht=2.25ex,dp=1ex,right]{date in head/foot}%
    \insertshortdate{}\hspace*{2em}
    \insertframenumber{} / \inserttotalframenumber\hspace*{2ex}
  \end{beamercolorbox}}%
  \vskip0pt%
}

% Configuración de Listings de Python para visualización pedagógica de código
\lstset{
    language=Python,
    basicstyle=\ttfamily\footnotesize,
    keywordstyle=\color{TecAzulOscuro}\bfseries,
    stringstyle=\color{TecRojo},
    commentstyle=\color{TecGris}\itshape,
    showstringspaces=false,
    breaklines=true,
    frame=lines,
    rulecolor=\color{TecAzulOscuro!30},
    backgroundcolor=\color{TecFondoBloque},
    aboveskip=0.5em,
    belowskip=0.5em,
    literate={á}{{\'a}}1 {é}{{\'e}}1 {í}{{\'i}}1 {ó}{{\'o}}1 {ú}{{\'u}}1
             {Á}{{\'A}}1 {É}{{\'E}}1 {Í}{{\'I}}1 {Ó}{{\'O}}1 {Ú}{{\'U}}1
             {ñ}{{\~n}}1 {Ñ}{{\~N}}1 {¿}{{?`}}1 {¡}{{!`}}1
}

% Metadatos institucionales por defecto
\title[Modelación Estadística]{Modelación Estadística para Ciencia de Datos e Ingeniería}
\author[J. Castillo Colmenares]{Juliho David Castillo Colmenares}
\institute[Tec de Monterrey]{Tecnológico de Monterrey \\ Escuela de Ingeniería y Ciencias}
\date{\today}

% Comandos y macros estadísticas compartidas para las presentaciones Beamer (Metropolis)

% Conjuntos numéricos básicos
\newcommand{\R}{\mathbb{R}}
\newcommand{\N}{\mathbb{N}}
\newcommand{\Q}{\mathbb{Q}}
\newcommand{\Z}{\mathbb{Z}}
\newcommand{\set}[1]{\left\{ #1 \right\}}
\newcommand{\abs}[1]{\left|#1\right|}
\newcommand{\norm}[1]{\left\| #1 \right\|}

% Comandos estadísticos y probabilísticos del curso
\newcommand{\Var}{\operatorname{Var}}
\newcommand{\cov}{\operatorname{Cov}}
\newcommand{\corr}{\rho}
\newcommand{\comb}[2]{\begin{pmatrix} #1 \\ #2 \end{pmatrix}}
\newcommand{\s}{\sigma}
\newcommand{\particion}{\mathfrak{P}}
\newcommand{\card}[1]{\#\left( #1 \right)}
\newcommand{\seq}[3]{\set{#1}_{#2}^{#3}}
\newcommand{\E}{\mathbb{E}}
\newcommand{\Prob}{\mathbb{P}}

% Entornos pedagógicos y cajas en español académico natural
\newenvironment{discusion}{%
  \begin{alertblock}{Pregunta de Discusión}%
}{%
  \end{alertblock}%
}

\newenvironment{reflexion}{%
  \begin{alertblock}{Reflexión para el Aula}%
}{%
  \end{alertblock}%
}

\newenvironment{paradoja}{%
  \begin{alertblock}{Paradoja de Exploración}%
}{%
  \end{alertblock}%
}

\newenvironment{motivacion}{%
  \begin{exampleblock}{Motivación: Ciencia de Datos en la Práctica}%
}{%
  \end{exampleblock}%
}

\newenvironment{retoclase}{%
  \begin{block}{Reto Rápido en Equipo}%
}{%
  \end{block}%
}


\title[Bondad de Ajuste $\chi^2$]{Sección 07.03: Pruebas de Bondad de Ajuste $\chi^2$}
\subtitle{Uniformidad, Parámetros Estimados y la Regla de Cochran}
\author[J. Castillo Colmenares]{Juliho Castillo Colmenares}
\institute{Tecnológico de Monterrey}
\date{\vspace{-1.2cm}}

\begin{document}

% Slide 1: Portada
\begin{frame}[plain]
  \titlepage
\end{frame}

% Slide 2: Hoja de Ruta
\begin{frame}{Hoja de Ruta --- Capítulo 07: Docimasia (Pruebas de Hipótesis)}
  \footnotesize
  ¿Dónde nos encontramos en el desarrollo modular del capítulo? \pause
  \begin{itemize}\itemsep=0.08em
    \item \textbf{07.01} Fundamentos: $H_0$ vs. $H_1$, Errores Tipo I y II, y Potencia \pause
    \item \textbf{07.02} Pruebas $Z$ y $t$ para Medias de Una y Dos Muestras \pause
    \item \textbf{07.03} Pruebas de Bondad de Ajuste $\chi^2$ \emph{(Hoy)} \pause
    \item \textbf{07.04} Tablas de Contingencia y Pruebas de Independencia
  \end{itemize}
  \vspace{-0.05cm}
  \begin{block}{Objetivo de la Sesión}
    Aprender a contrastar si un conjunto \emph{completo} de datos categóricos se ajusta a una distribución teórica propuesta --- uniforme, o con parámetros estimados de la propia muestra --- usando el estadístico $\chi^2$ de Pearson, respetando las condiciones de validez de la Regla de Cochran.
  \end{block}
\end{frame}

% Slide 3: Motivación
\begin{frame}{De un Parámetro a una Distribución Completa}
  \footnotesize
  \begin{motivacion}
    Hasta ahora hemos probado hipótesis sobre un solo parámetro (una media). Muchas preguntas de ciencia de datos son más ambiciosas: ¿siguen mis datos categóricos una distribución uniforme? ¿es razonable modelar estos conteos como Poisson? La prueba de bondad de ajuste $\chi^2$ contrasta el \emph{modelo completo}, no solo un parámetro aislado.
  \end{motivacion}

  \vspace{0.15cm}
  \pause
  \begin{itemize}\itemsep=0.1cm
    \item \textbf{Idea central:} comparar frecuencias \emph{observadas} contra las \emph{esperadas} bajo $H_0$. \pause
    \item \textbf{Aplicaciones:} justicia de un dado, ajuste a un modelo teórico (Poisson, Binomial), verificación de supuestos distribucionales.
  \end{itemize}
\end{frame}

% Slide 4: Estadístico y distribución
\begin{frame}{El Estadístico $\chi^2$ y su Distribución de Referencia}
  \footnotesize
  \begin{block}{Estadístico de Pearson}
    $$\chi^2 = \sum_{i=1}^k \frac{(O_i-E_i)^2}{E_i}, \qquad E_i = Np_i$$
  \end{block}
  \pause

  \vspace{0.1cm}
  \begin{alertblock}{Prueba Formal de Bondad de Ajuste}
    Bajo $H_0$, el estadístico converge asintóticamente a $\chi^2_\nu$ con $\nu=k-1-m$ grados de libertad, donde $m$ es el número de parámetros del modelo teórico \emph{estimados de la misma muestra}. Si $p_i$ son completamente conocidas de antemano, $m=0$ y $\nu=k-1$.
  \end{alertblock}
\end{frame}

% Slide 5: Regla de Cochran
\begin{frame}{La Regla de Cochran: Validez de la Aproximación}
  \footnotesize
  \begin{block}{Condición de Regularidad}
    La aproximación continua $\chi^2_\nu$ de la distribución discreta exacta del estadístico solo es confiable cuando las frecuencias esperadas no son demasiado pequeñas.
  \end{block}
  \pause

  \vspace{0.1cm}
  \begin{alertblock}{Regla Práctica de Cochran}
    \begin{itemize}\itemsep=0.05cm
      \item Ninguna $E_i$ debe ser menor a $1$.
      \item No más del $20\%$ de las celdas deben tener $E_i<5$.
    \end{itemize}
    Si se viola, la solución estándar es \textbf{fusionar categorías adyacentes} de baja frecuencia esperada hasta cumplir la regla, reduciendo correspondientemente los grados de libertad.
  \end{alertblock}
\end{frame}

% Slide 6: Casos particulares
\begin{frame}{Dos Escenarios: $H_0$ Simple vs. $H_0$ con Parámetro Estimado}
  \footnotesize
  \begin{block}{$H_0$ Simple (Probabilidades Conocidas)}
    Ejemplo: ¿es un dado justo? $H_0: p_i=1/6$ para $i=1,\ldots,6$. Aquí $m=0$, $\nu=k-1=5$.
  \end{block}
  \pause

  \vspace{0.1cm}
  \begin{alertblock}{$H_0$ con Parámetro Estimado}
    Ejemplo: ¿siguen los conteos una Poisson? Se estima $\hat\lambda=\bar X$ de la misma muestra antes de calcular las frecuencias esperadas. Aquí $m=1$, $\nu=k-1-1=k-2$: \textbf{cada parámetro estimado cuesta un grado de libertad adicional}.
  \end{alertblock}
\end{frame}

% Slide 7: Aplicaciones
\begin{frame}{Aplicaciones en Ciencia de Datos e Ingeniería}
  \footnotesize
  \begin{itemize}\itemsep=0.1cm
    \item \textbf{Control de calidad:} verificar si defectos por lote siguen una distribución esperada. \pause
    \item \textbf{Validación de modelos generativos:} ¿el ruido simulado sigue realmente la distribución supuesta? \pause
    \item \textbf{A/B/n testing:} verificar si el tráfico se distribuyó según las proporciones de asignación planeadas. \pause
    \item \textbf{Detección de fraude:} la Ley de Benford es, en esencia, una prueba de bondad de ajuste sobre el primer dígito.
  \end{itemize}
\end{frame}

% Slide 8: Lab Python (Bloque 1)
\begin{frame}[fragile]{Laboratorio en Python: Bondad de Ajuste Uniforme}
  \scriptsize
  Prueba de justicia de un dado, verificada contra \texttt{scipy.stats.chisquare} (\texttt{07.03\_goodness\_of\_fit\_tests.py}):
  {\lstset{basicstyle=\fontsize{4pt}{4.8pt}\selectfont\ttfamily,aboveskip=0.1em,belowskip=0.1em}
  \lstinputlisting[firstline=17, lastline=31]{../../code/07_pruebas_hipotesis/07.03_goodness_of_fit_tests.py}}
\end{frame}

% Slide 9: Lab Python (Bloque 2)
\begin{frame}[fragile]{Laboratorio en Python: Ajuste Poisson con Parámetro Estimado}
  \scriptsize
  Estimación de $\hat\lambda$ y ajuste de grados de libertad ($\nu=k-1-m$):
  {\lstset{basicstyle=\fontsize{3.6pt}{4.3pt}\selectfont\ttfamily,aboveskip=0.05em,belowskip=0.05em}
  \lstinputlisting[firstline=38, lastline=63]{../../code/07_pruebas_hipotesis/07.03_goodness_of_fit_tests.py}}
\end{frame}

% Slide 10: Lab Python (Bloque 3)
\begin{frame}[fragile]{Laboratorio en Python: Regla de Cochran y Fusión de Celdas}
  \scriptsize
  Detección de la violación de Cochran y corrección mediante fusión de categorías:
  {\lstset{basicstyle=\fontsize{4pt}{4.8pt}\selectfont\ttfamily,aboveskip=0.1em,belowskip=0.1em}
  \lstinputlisting[firstline=66, lastline=85]{../../code/07_pruebas_hipotesis/07.03_goodness_of_fit_tests.py}}
\end{frame}

% Slide 11: Salida Terminal
\begin{frame}[fragile]{Laboratorio en Python: Salida en Terminal}
  \scriptsize
  Salida estándar generada al ejecutar el script del laboratorio en Python:
  \vspace{0.02cm}
  \begin{block}{Salida Estándar en Consola}
    \fontsize{4pt}{5pt}\selectfont
    \begin{verbatim}
=== Block 1: GOF Against Uniform Hypothesis ===
chi2 statistic = 1.1000, df=5, p-value = 0.9541

=== Block 2: GOF With Estimated Parameter (Poisson) ===
lambda_hat = 0.9400, chi2 = 0.6977, df = 4-1-1 = 2, p-value = 0.7055

=== Block 3: Cochran's Rule and Cell Merging ===
Cells with E_i<5: 2 of 4 (50%) -- Cochran's Rule violated? True
Merged chi2 = 0.6890, df=2, p-value = 0.7086
    \end{verbatim}
  \end{block}
\end{frame}

% Slide 12A: Ejercicio Nivel 1 (Enunciado)

% Slide 12B: Ejercicio Nivel 1 (Resolución)

% Slide 13A: Ejercicio Nivel 2 (Enunciado)

% Slide 13B: Ejercicio Nivel 2 (Resolución)

% Slide 14A: Ejercicio Nivel 3 (Enunciado)
\begin{frame}{Ejercicio en Clase --- Nivel 3: Analítico (1/2)}
  \footnotesize
  \begin{block}{Problema --- Poisson con Parámetro Estimado (Enunciado, Sección 3.9)}
    Un ingeniero registra interrupciones diarias en $n=100$ días: $\{36,40,18,6\}$ para $k=0,1,2,\ge3$. Estime $\hat\lambda$, calcule las frecuencias esperadas bajo Poisson, y determine $\nu$ correctamente.
  \end{block}

  \vspace{0.1cm}
  \pause
  \begin{alertblock}{Estrategia}
    \scriptsize
    $\hat\lambda=\bar X$; los grados de libertad correctos son $\nu=k-1-m$ con $m=1$.
  \end{alertblock}
\end{frame}

% Slide 14B: Ejercicio Nivel 3 (Resolución)
\begin{frame}{Ejercicio en Clase --- Nivel 3: Analítico (2/2)}
  \scriptsize
  Estimamos $\hat\lambda$ y calculamos el estadístico ajustado: \pause

  \vspace{0.05cm}
  \begin{block}{Cálculo}
    \scriptsize
    $$\hat\lambda = \frac{0(36)+1(40)+2(18)+3(6)}{100} = 0.94, \qquad \chi^2 \approx 0.697, \qquad \nu = 4-1-1 = 2$$
  \end{block}

  \vspace{0.05cm}
  \pause
  \begin{alertblock}{Interpretación}
    \scriptsize
    Como $\chi^2=0.697 \ll \chi^2_{0.05,2}=5.991$, \textbf{se acepta con alta fidelidad} que las interrupciones siguen un proceso de Poisson con $\lambda=0.94$. Ignorar el ajuste $m=1$ habría inflado artificialmente $\nu$ a 3.
  \end{alertblock}
\end{frame}

% Slide 15A: Ejercicio Nivel 4 (Enunciado)
\begin{frame}{Ejercicio en Clase --- Nivel 4: Desafiante (1/2)}
  \footnotesize
  \begin{block}{Problema --- Fundamentación de $\nu=k-1$ (Enunciado, Sección 3.9)}
    Sea $\mathbf O=(O_1,\ldots,O_k)\sim\text{Multinomial}(N,\mathbf p)$. Justifique analíticamente por qué el estadístico de Pearson $Q=\sum Z_i^2$ se distribuye como $\chi^2_{k-1}$, y no $\chi^2_k$, explicando el rol de la ligadura $\sum(O_i-Np_i)=0$.
  \end{block}

  \vspace{0.1cm}
  \pause
  \begin{alertblock}{Estrategia}
    \scriptsize
    Observe que la suma fija $N$ impone una restricción lineal exacta sobre las $k$ desviaciones.
  \end{alertblock}
\end{frame}

% Slide 15B: Ejercicio Nivel 4 (Resolución)
\begin{frame}{Ejercicio en Clase --- Nivel 4: Desafiante (2/2)}
  \scriptsize
  Identificamos la ligadura lineal y su consecuencia geométrica: \pause

  \vspace{0.05cm}
  \begin{block}{Argumento}
    \scriptsize
    $\sum_i O_i = N \implies \sum_i(O_i-Np_i)=0$: el vector de desviaciones $\mathbf Z$ vive en un hiperplano de dimensión $k-1$, no en $\mathbb R^k$ completo. La última celda queda determinada por las primeras $k-1$.
  \end{block}

  \vspace{0.05cm}
  \pause
  \begin{alertblock}{Interpretación}
    \scriptsize
    Al proyectar sobre este hiperplano, el rango espectral efectivo de la forma cuadrática es $k-1$, no $k$. Por eso el estadístico de Pearson pierde exactamente un grado de libertad respecto a la intuición ingenua de sumar $k$ normales cuadráticas independientes.
  \end{alertblock}
\end{frame}

% Slide 16: Cierre
\begin{frame}{Cierre de Sección: Bondad de Ajuste $\chi^2$}
  \begin{reflexion}
    La prueba de bondad de ajuste contrasta un modelo distribucional \emph{completo}, no un parámetro aislado. Su validez depende críticamente de que las frecuencias esperadas no sean demasiado pequeñas (Regla de Cochran), y sus grados de libertad dependen exactamente de cuántos parámetros del modelo se estimaron de la misma muestra.
  \end{reflexion}

  \vspace{0.2cm}
  \pause
  \begin{block}{Lecciones Clave}
    \begin{itemize}\itemsep=0.08cm
      \item \textbf{Estadístico:} $\chi^2=\sum(O_i-E_i)^2/E_i$, con $\nu=k-1-m$.
      \item \textbf{Regla de Cochran:} ningún $E_i<1$; máximo 20\% de celdas con $E_i<5$; fusione si se viola.
      \item \textbf{Parámetros estimados:} cada uno estimado de la muestra resta un grado de libertad adicional.
    \end{itemize}
  \end{block}
\end{frame}

% Slide 17: Perspectiva Modular
\begin{frame}{Perspectiva Modular --- El Próximo Reto}
  \scriptsize
  \begin{retoclase}
    Ya sabemos contrastar si \emph{una} variable categórica sigue una distribución teórica. La Sección 07.04 cierra el Capítulo 07 extendiendo el mismo estadístico $\chi^2$ a \textbf{dos} variables categóricas simultáneas: las pruebas de independencia y de homogeneidad en tablas de contingencia.
  \end{retoclase}

  \vspace{0.08cm}
  \pause
  \begin{alertblock}{Preparación Recomendada}
    \begin{itemize}\itemsep=0.06cm
      \item Repasa el cálculo de frecuencias esperadas $E_{ij}=R_iC_j/N$ en tablas bidimensionales.
      \item Reflexiona sobre la diferencia conceptual entre diseñar el muestreo por filas o por columnas fijas.
    \end{itemize}
  \end{alertblock}
\end{frame}

\end{document}
